\chapter{Network Stability and Model Selection}

\begin{enumerate}
    \item Previous chapter focused on how to estimate and interpret moderated network models. But many important questions remain for the researcher hoping to apply such an analysis to their data:
    \begin{enumerate}
        \item Among all possible model structures to choose from, what is the optimal way to decide which model to interpret? That is, on what grounds can we engage in model selection?
        \item Given a model that has been chosen, how can we assess the accuracy of the parameter estimates? More specifically, how can we assess the degree to which our model is subject to sampling variation? (Accuracy = proneness to sampling variation)
        \item How stable is a particular model structure, and how likely is it to replicate in another sample? (Stability -- robust structure in smaller samples)
        \item Finally, how can we address issues of statistical power and determine necessary sample sizes a priori?
    \end{enumerate}
    \item Each of these questions will be dealt with in this chapter and accompanied by simulations to demonstrate various ways of exploring a dataset, designing confirmatory network-based studies, and providing a comprehensive analysis of moderated networks that address common issues/questions in modeling complex networks using real data.
\end{enumerate}

\section{Model Selection}
First step is talking about the importance of selecting a good model. The key is to balance parsimony with predictive accuracy, particularly when this is an out-of-sample based form of accuracy. Researchers who work with psychological networks often make an assumption that the `true' network or `true' model is sparse. This means that researchers typically anticipate that only a moderate to small proportion of the total number of possible relationships in a given network actually map onto real causal connections, rather than spurious associations. The implication of this assumption is that researchers are faced with the often difficult task of determining which model to choose and interpret.

Model selection is certainly not a new statistical problem by any means. In fact it is standard procedure in most situations of modeling multivariate or complex phenomena. Common techniques that have been used for decades include forward, backward, and best-subset selection, all of which involve searching a large probability space and comparing a multitude of models to find which fits some pre-defined criterion (e.g., to minimize the sum of squared errors, or some other information criterion). Overall, we can distinguish between 3 general approaches to model selection: (a) Thresholding (pruning), (b) Model search, and (c) Regularization.

\subsection{Thresholding}
Thresholding is probably the simplest and most widely used form of model selection. The most basic example of this would be when a researcher fits a regression model with multiple predictors, but then only interprets the effect estimates which, under the null hypothesis, are sufficiently low in probability (e.g., $p < .05$). The most common way of implementing this method is by simply restricting one's interpretation of the model to its significant parameters (the threshold of which will be determined by the researcher). Alternatively, one might assess which parameters of a given model are significant and then re-fit the model while only including those parameters. In highly complex models with many variables, however, this routine can be implemented recursively until the researcher cannot add or subtract any parameters on the basis of the pre-defined threshold.

\subsection{Model Search}
Thresholding is a rather crude way of performing model selection, as the researcher commits to making binary decisions about model parameters based on a fixed significance threshold, and therefore does not think cautiously about the effects of including/excluding any given parameter, and whether certain combinations of included parameters lead to unexpectedly well-fitting models. The central question driving \textit{model search} as a model selection process  is: what model explains the data appropriately while also being as simple as possible?

The primary way of answering this question is by selecting an information criterion that reflects some degree of trade-off between parameter accuracy and parsimony, and then evaluating a variety of models on the basis of that criterion to find the one that satisfies some objective. The challenge here, however, becomes the model space (i.e., the space of all possible models which are being searched). For an undirected network model, for instance, there are $2^{p(p - 1)/2}$ possible models given a network with $p$ nodes. This is another common approach, but can become intractable quickly. Procedures such as backward and forward selection arose to cope with (not solve) this issue. 

\subsection{Regularization}
This is perhaps the method that is least frequently used within the psychological sciences, however it is the method that is extremely prominent within the network-modeling community. The downside of thresholding and model search is that estimation and selection are performed in separate steps. First a model is estimated, and then parameters are evaluated as candidates for inclusion/exclusion in the model, and then a new model is fit. With regularization, the idea is to use a different estimator---specifically a \textit{penalized} estimator---such that estimation and variable selection are performed in the same step. This can be achieved via that LASSO, which stands for Least Absolute Shrinkage and Selection Operator. The LASSO is extremely common within the psychological network literature, and is almost a default in the current community.

The upsides of the LASSO is that it is highly efficient and performs automatic variable selection when parameters for a model are estimated. The downsides, however, are that (a) the estimates from the LASSO are fundamentally biased (e.g., all estimates are shrunk closer to, or exactly to, zero), and (b) this prevents proper standard errors from being determined for the estimates; the sampling distribution of a LASSO'd parameter is unknown, and cannot be approximated through bootstrapped resampling. 

\subsection{The Problem with Moderated Networks}
Although these approaches for selecting a model among an array of potential candidates are widely studied and applied to many different problems. However, in the case of higher-order interaction terms (whether these are multiplicative interactions between multiple variables, or polynomial effects e.g., with respect to time), these semi-automated approaches to model selection can be problematic. Specifically, it is widely acknowledged that omitting constituent main effects associated with a multiplicative interaction term when it is included in a model can lead to biased coefficients of the relationships between the predictors and the outcome(s). [COULD BE COOL TO HAVE AN EXAMPLE HERE SHOWING HOW OMITTING MAIN EFFECTS LEADS TO BIASED ESTIMATES].

\subsection{Summary}
Basically, (a) model selection is important, (b) many methods are available, but most of these fall short when it comes to evaluating models with higher-order interactions. What we need are algorithms and methods that enforce a structural hierarchy on the model search and/or regularization process. So, now to introduce the hierarchical LASSO!

\section{Variable Selection with the Hierarchical LASSO}
First, a discussion of the hierarchical LASSO, group LASSO, etc. Then, a deep dive into the algorithms:
\begin{enumerate}
    \item Multi-sample splitting with p-value adjustment
    \begin{enumerate}
        \item Split sample in half. 
        \item Fit hierarchical LASSO on one half
        \item Then, take selected variables from first half and fit unregularized model with those parameters to the second half.
        \item Save estimates from each iteration, but particularly the p-values. If a variable is not selected in a given iteration, set the p-value to 1. 
        \item Finally, use a quantile-based adjustment to the resulting distribution of p-values associated with each parameter in order to control the false discovery rate.
        \item Adjusted confidence intervals can be computed on a similar basis; the final model is then chosen on the basis of either the adjusted p-values or confidence intervals.
    \end{enumerate}
    \item Bootstrapping with p-value adjustment
    \begin{enumerate}
        \item Same shit as above but with bootstrapping... yippee. Although it is interesting to think about what the differences are.. in theory, the subsampling approach should lead to better out-of-sample prediction accuracy, while on the other hand the bootstrapping method would have more power and so may actually generate more stable estimates. [EMPIRICAL QUESTION]
    \end{enumerate}
    \item Stability Selection
    \begin{enumerate}
        \item My favorite version -- simultaneous stability selection. 
        \item This involves undertaking the same multi-sample splitting approach (subsampling), but a separate hierarchical LASSO model is fit to each subsample of the data. Then, for each parameter we assign a 0 or a 1 based on whether or not that variable was selected in both halves of the data.
        \item Once multiple iterations are completed, for each variable we can take the proportion of times that the parameter was chosen in each half of the data across iterations.
        \item Some threshold is determined here... some way to control FDR, but not clear on exactly how to do that.
    \end{enumerate}
\end{enumerate}

It's worth adding that these all act as model selection procedures or algorithms, but in addition to this we must have some criterion on which to base our decisions in the search procedure. This is likely going to be minimizing an information criterion, or minimizing cross-validated error.

%\chapter{Variable Selection in Moderated Networks}
One key assumption in psychological network modeling is that the `true' network has a \textit{sparse} structure. There are two primary reasons for this assumption: (a) First, in many circumstances it seems most likely that there are only a few important (out of many possible) causal pathways among the variables under study, and (b) second, sparse models are much easier to interpret than dense models. The latter point is truly a pragmatic consideration, but is nevertheless important because of the rapidly-increasing degree of complexity in network models as we add new variables to the system (especially when higher-order interactions are included). We need to be able to `hone in' on the key relationships and focus on the primary drivers of a system's behavior in order to better design confirmatory and experimental studies to explore it further. 

Assuming that psychological networks are sparse is often a starting point for network modeling in psychology. This usually seems like a reasonable assumption to make, and it may even be necessary due to the highly exploratory nature of this field of research. Given the large number of relationships that are often being estimated, techniques from machine learning and `big data' come into play more frequently, and thereby bring along new sets of considerations. Specifically, the parameters in network models are often estimated using penalized regression techniques---most notably, the \textit{LASSO}. This approach was designed for high-dimensional data analysis (where $p \gg n$), and offers some notable advantages over other approaches when the goal is to estimate some kind of sparse statistical structure. I will discuss this method in the following sections, as well as some of its drawbacks, and how it can be used in cases where we are interested in studying moderation effects.

\section{Variable Selection with the LASSO}
LASSO is an acronym that stands for `Least Absolute Shrinkage and Selection Operator', and is a method of penalized regression that offers a powerful alternative to ordinary least squares (OLS) \cite{tibshirani1996regression}. The key function of the LASSO is to increase prediction accuracy and model interpretability by shrinking coefficients of the predictors toward zero---importantly, this procedure will often shrink some coefficients to be \textit{exactly} zero, thereby performing automatic variable selection and creating a model that's easier to interpret. The general form of the LASSO can be stated as:
\begin{equation}
    \hat{\bbeta} = \argmin_{\bbeta} \lb \sum_{i=1}^n (y_i - \bx_i \bbeta)^2 + \lambda\sum_{j=1}^p |\beta_j| \rb,
\end{equation}
where $\hat{\bbeta}$ is the vector of coefficients that minimizes the residual sum of squares (RSS) along with an added penalty that weights the sum of the absolute values of the coefficients with a parameter $\lambda$. The objective represented by the formula above could thus be conceptualized as finding the coefficients that minimize: RSS + penalty.

The tuning parameter $\lambda$ governs the strength of the penalty, and thereby the degree of shrinkage that is applied to the coefficients. Larger values of $\lambda$ enforce a larger penalty, which is likely to translate into more parameters being set to zero. At the most extreme, all coefficients may be set to zero, leading to an intercept-only model, and conversely when this parameter is small (or is itself zero), the model will converge toward unpenalized OLS estimates. The optimal value of this parameter is often determined through $k$-fold cross-validation, wherein an array of $\lambda$ values are used in fitting models to different subsets of the data and then tested on $k$ hold-out sets to determine which value results in the lowest cross-validated mean-squared error. Information criteria (such as the AIC and BIC) may be used as an alternative to cross-validation for this purpose, but in all cases the goal is to find which value of $\lambda$ results in the model with the best balance between accuracy and parsimony. 

The primary attractive features of the LASSO are that it narrows down the number of predictors to only the most important (i.e., only those with the strongest relationships to the outcome) by setting the coefficients for some predictors to zero. Not only does this have the advantage of making models more interpretable, but it can also increase prediction accuracy by reducing the possibility of overfitting---that is, fitting a model that gives too much weight to the noise within some particular dataset, and thereby creates a model that fails to generalize well to new cases. 

In the case of modeling psychological networks, these attributes are desirable because of the large number of possible coefficients (`edges') that may be used to describe even a small network. In the most simple case where we are fitting a network model to cross-sectional data, for instance, given $p$ variables (or `nodes') there will always be $p(p - 1)/2$ edges being estimated. When there are only 3 nodes in the network, this means that there will only be 3 edges. But with 6 nodes we now have 15 edges, and with 10 nodes there will be 45 edges. This number grows exponentially as we add more nodes to the network, and is therefore something that researchers aim to minimize as much as possible in order to highlight only the most important relationships. 

\subsection{Disadvantages of the LASSO}
Although the LASSO has some very desirable properties for aiding in interpretation and increasing predictive accuracy, this comes at the cost of producing coefficient estimates that are biased towards zero, and as a result may not provide a good representation of their `true' values. While this is not necessarily a problem when prediction is the goal, it does create a major obstacle for inference. Specifically, there is no known way to determine the standard errors of estimates derived from the LASSO, and thus no way to perform significance tests that can be assumed to reflect the characteristics of the associated sampling distributions. Due to the unique bias of the LASSO estimates (i.e., that it can lead to some parameters being estimated at zero), even bootstrapped resampling fails to be a valid method for determining the distributions of the parameters. This method may be effective when the true coefficients are especially large, or truly equal to zero, but for cases in between it will usually generate asymmetric distributions that may not actually describe the population-level associations between the candidate variable(s) and outcome of interest. 

A related disadvantage of the LASSO, although not unique to the LASSO, is inconsistent model selection. In cases where small changes in the tuning parameter $\lambda$ lead to different variables being included or excluded from the final model, it may be difficult to make strong inferences about the value of those variables for understanding the phenomena under study. To my knowledge, no variable selection technique is immune to this concern, however, and it may be a persistent challenge for any approach of this nature (e.g., best subset selection). However, despite the fact that bootstrapping does not seem to solve the inference problem, it does provide the researcher with a sense of the selection procedure's stability with regard to each of the predictors. One possible way of increasing confidence in the chosen variables is by reporting the selection rates with reference to each predictor; for instance, we may have greater confidence in choosing variables that are selected for the best model within 90\% of subsamples as opposed to those that are chosen only 15\% of the time.


\subsection{Variable Selection and Interaction Terms}
With regard to the current research, it is a problem to employ automatic variable selection when higher-order terms (such as interactions) are included in a model, as they are treated indiscriminately and without regard to their associated lower-order terms. Specifically, this means that models may be selected wherein interaction terms are included without the inclusion of the related main effects. This is problematic, as psychologists are generally interested in determining whether interactions between variables are significant \textit{over and above} the main effects that constitute them. Thus, when interested in assessing the influence of moderator variables on some outcome of interest, it is important to employ a variable selection method that only chooses models which respect the structure of the predictor variables. 

One method that accomplishes this goal is the \textit{hierarchical LASSO}, which was specifically designed for addressing cases of model selection that include higher-order interactions or polynomial terms \cite{lim2015learning,bien2013lasso}. Available approaches that use this method can be specified to only choose models that obey a \textit{strong hierarchy}, meaning that whenever an interaction or higher-order term is selected for a given model, the associated lower-order term(s) are included as well. Two statistical packages that implement this in the R programming language include: \code{hierNet} \cite{bien2013lasso}, and \code{glinternet} \cite{lim2015learning}. Both of these packages utilize an overlapping group LASSO, wherein the LASSO is applied to predictor variables on the basis of their grouping into sets of main effects + interactions, or main effects and their associated polynomial terms. 

Both of these packages produce similar results in a variety of different settings, but \code{glinternet} exhibits a major computational advantage in terms of computing time. Furthermore, while both packages afford researchers to search the entire model space of all possible interactions among all predictor variables, \code{glinternet} also allows one to select a subset of variables (or a single variable) to serve as candidates for interactions. This is particularly useful when the researcher has specific variables in mind which may serve as moderators.

\section{Stability Selection}
Goal is to choose one element from a set of models
\begin{equation}
    \{\hat{S}^{\lambda};\lambda\in\Lambda\}
\end{equation}
where $\Lambda$ is the full set of regularization parameters considered. There might be two problems: (1) First the correct model $S$ might not even be in the set being searched. (2) Second, even if it is a member, it is not clear to know what the ideal values of $\lambda$ would be to detect $S$. With stability selection, we do not simply select one model. Instead the data are perturbed (by subsampling) many times and we choose all structures of variables that occur in a large fraction of the resulting selection sets. 

For a given cut-off $\pi_{\text{thr}}$ with $0 < \pi_{\text{thr}} < 1$ and a set of regularization parameters $\Lambda$, the set of stable variables is defined as 
\begin{equation}
    \hat{S}^{\text{stable}} = \{k:\max\limits_{\lambda\in\Lambda}(\hat{\Pi}_k^{\lambda})\} \geq \pi_{\text{thr}} 
\end{equation}
We keep variables with a high selection probability and disregard those with low selection probabilities. The exact cut-off $\pi_{\text{thr}}$ is a tuning parameter but results vary surprisingly little for sensible choices of the range. Nor do results seem to depend strongly on regularization. Barbieri and Berger (2004) showed predictive optimality results with the \textit{median probability model}, consisting of variables which have posterior probability of being in the model of $\frac{1}{2}$ or greater. This is different than choosing the model with the highest posterior probability. 

